\documentclass[letterpaper]{report}

\usepackage{tikz, background, titling, setspace}
\usepackage{tabu}
\usetikzlibrary{calc}
\usepackage{fancyhdr}
\usepackage{tcolorbox}
\usepackage{layout}
\usepackage[left=1in,right=0.5in,top=0.625in]{geometry}
\usepackage{pgf}
\usepackage{pgfpages}
\usepackage{pgfplots}
\usepackage{graphicx}
\usepackage{amsfonts}
\usepackage{amsmath}
\usepackage{amssymb}
\usepackage{tikz,pgfplots}
\usepackage[abspage,user,lastpage]{zref}
\usepackage{lastpage}
\usepackage{marginnote}
\usepackage{eso-pic}
\usepackage{ulem}
\usepackage{tikz}
\usepackage{circuitikz}
\usepackage{tikz-timing}
\usepackage{wrapfig}
\usepackage[export]{adjustbox}
\usepackage{listings}
\usepackage{xcolor}
\usepackage{soul}
\usepackage{cite}
\usepackage{url}
\usepackage[absolute]{textpos}
\usepackage{amssymb}  % provides Rbb, etc.
\usepackage{latexsym}  % provides \Box
\usepackage{array}
\usepackage{siunitx}
\usepackage{steinmetz}
\usepackage{mdframed}
\usepackage{changepage}


% Defined Colors to be used throughout
% Change the color by uncommenting the color you want and commenting out the one you don't
  % Page Background color
    % #FFEC8B (LightGoldenrod1)
        % \definecolor{PageColorz}{RGB}{255,236,139}
    % #FFFE4  (darker YELLOW) -Defualt-
      %\definecolor{PageColorz}{RGB}{255,255,200}
    % #FFFFDC (Light Yellow1) -Recomendable-
      \definecolor{PageColorz}{RGB}{255,255,220}
    % #FFFFED (Light Yellow)
      % \definecolor{PageColorz}{RGB}{255,255,237}
    % #fffff5 (Lighter Yellow)
      % \definecolor{PageColorz}{RGB}{255,255,245}
    % #ffffff (White)
      % \definecolor{PageColorz}{RGB}{255,255,255}
    % #BFEFFF (LightBlue1)
      % \definecolor{PageColorz}{RGB}{191,239,255}
    % #E0FFFF (LightCyan)
      % \definecolor{PageColorz}{RGB}{224,255,255}
    % #104E8B (DodgerBlue4)
	    % \definecolor{PageColorz}{RGB}{16,78,139}
    % #097969 (CadmiumGreen)
      % \definecolor{PageColorz}{HTML}{097969}
    % #AFE1AF (CeladonGreen) 
      % \definecolor{PageColorz}{HTML}{AFE1AF}

    % MEME colorz (PINK) #FFb6C1
      % \defineColor{PageColorz}{HTML}{FFB6C1}
	\pagecolor{PageColorz}

  % Line Colors 
    % #b94249 (Green) -Default-
      \definecolor{MyLines}{RGB}{73,185,66}
    % #6495ed (Blue)
      % \definecolor{MyLines}{RGB}{97,145,179}
    % #104E8B (DodgerBlue4)
      % \definecolor{MyLines}{RGB}{16,78,139}
    % #E0FFFF (LightCyan)
      % \definecolor{MyLines}{RGB}{224,255,255}
    % #000000 (Black)
      % \definecolor{MyLines}{RGB}{0,0,0}
    % #FFFFFF (White)
      % \definecolor{MyLines}{RGB}{255,255,255}

  % Default Text Color
    % #b94249 (Green)
      % \definecolor{MyColor}{RGB}{73,185,66}
    % #6495ed (Blue)
      % \definecolor{MyColor}{RGB}{97,145,179}
    % #000000 (Black) -Default-
      \definecolor{MyColor}{RGB}{0,0,0}
  
    \color{MyColor}

  % GridColor
    % #DFFFFF (UltraLightGrey) -Defualt-
      \definecolor{GridColor}{HTML}{DFFFFF}
    % #D3D3D3 (DarkishGrey)
      % \definecolor{GridColor}{HTML}{D3D3D3}

  % Extra Defined Color
  \definecolor{Extra}{RGB}{16,78,160}
  % \definecolor{Extra}{RGB}{205,245,255}

% A compatability option for pgfplots
\pgfplotsset{compat=1.18}

% Document Properties
\newcommand{\course}{ECE 3620}
\newcommand{\studentName}{Petersen, Kelson}
\newcommand{\aNumber}{A02346768}
\newcommand{\assnNumber}{10}
\newcommand{\w}{\omega}
\newcommand{\Ohm}{\Omega}
\newcommand{\D}{^\circ}
% Draw the Background lines
\backgroundsetup{%
 position=current page.center,
 angle=0,
 scale=1,
 contents={%
  \begin{tikzpicture}%
    [
      normal lines/.style={MyLines, thick},
      grid lines/.style={GridColor, very thin},
      every node/.append style={black, align=center, opacity=1}
    ]
    \foreach \y in {0.5, 1,...,27} {  \draw[grid lines] (2.5cm,\y cm) -- (20.5cm,\y cm); }
    \foreach \x in {2.5, 3,...,20.5}          {  \draw[grid lines] (\x cm,0cm) -- (\x cm,27 cm); }
	% Draw the background lines of the document
	\draw[normal lines] (3.166in,27cm) -- (3.166in,11in);
	\draw[normal lines] (5.333in,27cm) -- (5.333in,11in);
    \draw[normal lines] (2.5cm,0) -- (2.5cm,11in);
    \draw[normal lines] (20.5cm,0) -- (20.5,11in);
    \draw[normal lines] (0in,27cm) -- (21.5cm,27cm);
	% Make Nodes for the header to appear
	\node (t) [font=\footnotesize, anchor=center] at (4.25in, 27.5cm) {\course; Hw \assnNumber};
    \node (t) [font=\footnotesize, anchor=center] at (5cm, 27.5cm) {November\hspace{1.75em}, 2025};
    \node (t) [font=\footnotesize, anchor=center] at (15cm, 27.5cm) {};
	\node (t) [font=\footnotesize, anchor=center] at (21cm, 27.5cm) {};
  \end{tikzpicture}%
}}


% Pagestyle for no header 
\pagestyle{empty}

% Spacing Controls
\setlength{\parindent}{0pt}
% New environments

\newcounter{ProblemNumber}[section]
\newenvironment{problem}[1][]{
  \par\noindent\textbf{Problem #1 \refstepcounter{ProblemNumber}:}\par\smallskip
  \begin{adjustwidth}{0.15cm}{0.15cm}
}{
  \end{adjustwidth}
  \vspace{6pt}
  \noindent\textcolor{MyLines}{\rule{\textwidth}{0.6pt}}
  \medskip
}
%\newcounter{SolutionNumber}[section]
\newenvironment{answer}[1][]{
  \setlength{\leftskip}{0.125in}
  \setlength{\rightskip}{0.125in}
  %\textbf{Solution #1 \refstepcounter{ProblemNumber}:} \\ 
}
{
  \newline
  \textcolor{MyLines}{\hrule}
  \vspace{2pt}
  \textcolor{MyLines}{\hrule}
  \vspace{6pt}
}

\begin{document}
\begin{problem}[{5.19}]
	\textbf{Given:} The figures below\\
  \begin{minipage}{0.5\textwidth}
    \includegraphics[scale=0.125]{fig5.6.png}
  \end{minipage}
  \begin{minipage}{0.5\textwidth}
    \includegraphics[scale=0.125]{fig5.7.png}
  \end{minipage}
  \begin{minipage}{0.5\textwidth}
    \includegraphics[scale=0.125]{fig5.8.png}
  \end{minipage}
  \begin{minipage}{0.5\textwidth}
    \includegraphics[scale=0.125]{fig5.9.png}
  \end{minipage}
  \begin{minipage}{0.5\textwidth}
    \includegraphics[scale=0.125]{fig5.10.png}
  \end{minipage}
  \begin{minipage}{0.5\textwidth}
    \includegraphics[scale=0.125]{fig5.11.png}
  \end{minipage}
	\textbf{Find:} Which of the waveforms above will exhibit the Gibbs oscillation phenomenon when represented by a Fourier series? Why?
\end{problem}
% \begin{answer}
% 	Answer Section
% \end{answer}
\newpage
\begin{problem}[{5.39}]
	\textbf{Given:} $A=5$ and $T=3$ s.
  \begin{center}
    \includegraphics[scale=0.175]{fig5.39.png}
  \end{center}
	\textbf{Find:} Determine the Fourier transform.
\end{problem}
% \begin{answer}
% 	Answer Section
% \end{answer}
\newpage
\begin{problem}[{5.40}]
	\textbf{Given:} $A=10$ and $T=6$ s.
  \begin{center}
    \includegraphics[scale=0.175]{fig5.40.png}
  \end{center}
	\textbf{Find:}Determine the Fourier transform.
\end{problem}
% \begin{answer}
% 	Answer Section
% \end{answer}
\newpage
\begin{problem}[{5.42}]
	\textbf{Given:} $A=2$ and $T=12$ s.
  \begin{center}
    \includegraphics[scale=0.175]{fig5.42.png}
  \end{center}
	\textbf{Find:}Determine the Fourier transform.
\end{problem}
% \begin{answer}
% 	Answer Section
% \end{answer}
\newpage
\begin{problem}[{5.46}]
  \textbf{Given:} $A=5$, $T=1$ s, and $\alpha=10~s^{-1}$.
  \begin{center}
    \includegraphics[scale=0.175]{fig5.46.png}
  \end{center}
	\textbf{Find:}Determine the Fourier transform.
\end{problem}
% \begin{answer}
% 	Answer Section
% \end{answer}
\newpage
\begin{problem}[{5.47}]
	\textbf{Given:} $A=10$ and $T=2$ s.
  \begin{center}
    \includegraphics[scale=0.175]{fig5.47.png}
  \end{center}
	\textbf{Find:}Determine the Fourier transform.
\end{problem}
% \begin{answer}
% 	Answer Section
% \end{answer}
\newpage
\begin{problem}[{5.50}]
	\textbf{Given:}
  \begin{itemize}
    \item[a.] $f(t) = e^{-\alpha t}\sin(\omega_1t)\cos(\omega_2t)u(t)$
    \item[b.] $g(t) = te^{-\alpha t},~~0\leq t \leq 10\alpha$
  \end{itemize}
  \textbf{Find:} the Fourier transform of the following signals with $\alpha=0.5~s^{-1}$, $\omega_1=4\text{rad}/s$, and $\omega_2=2\text{rad}/s$
\end{problem}
% \begin{answer}
% 	Answer Section
% \end{answer}
\newpage
\begin{problem}[{5.66}]
	\textbf{Given:}
  \begin{align}
    \int^\infty_{-\infty}\frac{{\sin}^2(at)}{{(\pi t)}^2}dt = \frac{a}{\pi}
  \end{align}
	\textbf{Find:} Show the above is true if $a > 0$.
\end{problem}
% \begin{answer}
% 	Answer Section
% \end{answer}

\end{document}
